\thispagestyle{empty}
\begin{center}
    \textbf{\fontsize{14pt}{16pt}\selectfont ĐÁNH GIÁ QUYỂN ĐỒ ÁN TỐT NGHIỆP}\\
    \textit{\fontsize{12pt}{14pt}\selectfont (Dùng cho cán bộ phản biện)}
\end{center}

\vspace{10pt}
\noindent Tên giảng viên đánh giá: \dotfill \\
Họ và tên sinh viên: \dotfill\ MSSV: \makebox[2.5cm]{\dotfill} \\
Tên đồ án: \dotfill \\
\makebox[\textwidth]{\dotfill}

\vspace{10pt}
\noindent \textbf{Chọn các mức điểm phù hợp cho sinh viên trình bày theo các tiêu chí dưới đây:}\\
\textit{Rất kém (1); Kém(2); Đạt(3); Giỏi(4); Xuất sắc(5)}

\vspace{5pt}
\begin{center}
\fontsize{9pt}{11pt}\selectfont
\renewcommand{\arraystretch}{1.2}
\begin{tabularx}{\textwidth}{|c|X|c|c|c|c|c|}
\hline
\multicolumn{7}{|l|}{\textbf{Có sự kết hợp giữa lý thuyết và thực hành (20)}} \\ \hline
1 & Nêu rõ tính cấp thiết và quan trọng của đề tài, các vấn đề và các giả thuyết (bao gồm mục đích và tính phù hợp) cũng như phạm vi ứng dụng của đồ án & 1 & 2 & 3 & 4 & 5 \\ \hline
2 & Cập nhật kết quả nghiên cứu gần đây nhất (trong nước/quốc tế) & 1 & 2 & 3 & 4 & 5 \\ \hline
3 & Nêu rõ và chi tiết phương pháp nghiên cứu/giải quyết vấn đề & 1 & 2 & 3 & 4 & 5 \\ \hline
4 & Có kết quả mô phỏng/thực nghiệm và trình bày rõ ràng kết quả đạt được & 1 & 2 & 3 & 4 & 5 \\ \hline
\multicolumn{7}{|l|}{\textbf{Có khả năng phân tích và đánh giá kết quả (15)}} \\ \hline
5 & Kế hoạch làm việc rõ ràng bao gồm mục tiêu và phương pháp thực hiện dựa trên kết quả nghiên cứu lý thuyết một cách có hệ thống & 1 & 2 & 3 & 4 & 5 \\ \hline
6 & Kết quả được trình bày một cách logic và dễ hiểu, tất cả kết quả đều được phân tích và đánh giá thỏa đáng & 1 & 2 & 3 & 4 & 5 \\ \hline
7 & Trong phần kết luận, tác giả chỉ rõ sự khác biệt (nếu có) giữa kết quả đạt được và mục tiêu ban đầu đề ra đồng thời cung cấp lập luận để đề xuất hướng giải quyết có thể thực hiện trong tương lai & 1 & 2 & 3 & 4 & 5 \\ \hline
\multicolumn{7}{|l|}{\textbf{Kỹ năng viết quyển đồ án (10)}} \\ \hline
8 & Đồ án trình bày đúng mẫu quy định với cấu trúc các chương logic và đẹp mắt (bảng biểu, hình ảnh rõ ràng, có tiêu đề, được đánh số thứ tự và được giải thích hay để cập đến; căn lề thống nhất, có dấu cách sau dấu chấm, dấu phẩy v.v.), có mở đầu chương và kết luận chương, có liệt kê tài liệu tham khảo và có trích dẫn đúng quy định & 1 & 2 & 3 & 4 & 5 \\ \hline
9 & Kỹ năng viết xuất sắc (cấu trúc câu chuẩn, văn phong khoa học, lập luận logic và có cơ sở, từ vựng sử dụng phù hợp v.v.) & 1 & 2 & 3 & 4 & 5 \\ \hline
\multicolumn{7}{|l|}{\textbf{Thành tựu nghiên cứu khoa học (5) \textit{(chọn 1 trong 3 trường hợp)}}} \\ \hline
10a & Có bài báo khoa học được đăng hoặc chấp nhận đăng/Đạt giải SVNCKH giải 3 cấp Viện trở lên/Có giải thưởng khoa học (quốc tế hoặc trong nước) từ giải 3 trở lên/Có đăng ký bằng phát minh, sáng chế & \multicolumn{5}{c|}{5} \\ \hline
10b & Được báo cáo tại hội đồng cấp Viện trong hội nghị SVNCKH nhưng không đạt giải từ giải 3 trở lên/Đạt giải khuyến khích trong các kỳ thi quốc gia và quốc tế khác về chuyên ngành (VD: TI contest) & \multicolumn{5}{c|}{2} \\ \hline
10c & Không có thành tích về nghiên cứu khoa học & \multicolumn{5}{c|}{0} \\ \hline
\multicolumn{2}{|l|}{\textbf{Điểm tổng}} & \multicolumn{5}{r|}{/50} \\ \hline
\multicolumn{2}{|l|}{\textbf{Điểm tổng quy đổi về thang 10}} & \multicolumn{5}{r|}{} \\ \hline
\end{tabularx}
\end{center}

\newpage
\thispagestyle{empty}

\vspace*{10pt}
\noindent \textbf{Nhận xét khác của cán bộ phản biện}\\
\noindent\rule{\textwidth}{0.1mm}\\
\noindent\rule{\textwidth}{0.1mm}\\
\noindent\rule{\textwidth}{0.1mm}\\
\noindent\rule{\textwidth}{0.1mm}\\
\noindent\rule{\textwidth}{0.1mm}

\vspace{25pt}
\hfill \parbox{5cm}{
    \centering
    Ngày: \dots / \dots / 20\dots \\
    \vspace{10pt}
    \textbf{Người nhận xét}\\
    \textit{\fontsize{10pt}{12pt}\selectfont (Ký và ghi rõ họ tên)}
}
\newpage
